\chapter{Resolution, Precision, and the Continuum}
\label{chap:resolution-precision-continuum}

Here the discrete most directly confronts the smooth. The continuum --- the unbroken real line physics is
written on --- stands here as a shadow the discrete casts, never the substrate beneath it: an analogy, the most
useful one the instrument keeps and the most carefully fenced, offered as an alternative and never as a
disproof --- the frame ours to offer, the smooth limit the refinements approach the world's. Six effects make
the case: a resolution floor below which nothing varies, a precision that is reproducibility and not truth, a
derivative that needs no vanishing quantity, a continuum that is a limit of refinements rather than a ground, a
length that depends on the ruler, and an irrational diagonal no enumeration can finish.

\section{The Berkeley--Galileo Effect: a resolution floor}
\label{se:berkeley-galileo}

No finite instrument resolves arbitrarily small variation. An instrument of resolution $\epsilon$ records not
the state but its shadow,
\begin{equation}
  \text{shadow} = \left\lfloor \frac{\text{state}}{\epsilon} \right\rfloor,
\end{equation}
the state divided into whole cells of width $\epsilon$ and the remainder discarded --- the cell width ours, the
state it coarsens the world's; any variation smaller than
$\epsilon$ leaves the shadow unchanged, so below the threshold multiple candidate descriptions are
indistinguishable, separable only by a finer instrument. Galileo, grinding lenses to push the floor down, saw
moons around Jupiter the naked eye could not --- not because the moons were infinitesimal but because they lay
below the eye's $\epsilon$ and above the telescope's. And Berkeley, a century and more later, objected to Newton's
infinitesimals as ``ghosts of departed quantities'' --- magnitudes treated as nonzero to form a ratio and then
as zero to be rid of. The objection is answered here by a floor: there are no departed quantities, only
variations below resolution, and a finite instrument cannot see them --- not because they vanish, but because
they were never above the floor.

The honest floor is a finite count obligation --- the shadow is a count of $\epsilon$-cells, and
sub-$\epsilon$ variation adds no count --- the ruler ours to bring, what lies beneath its floor the world's.
The reading is that resolution is a floor, not a limit approached.
Below $\epsilon$ nothing varies that the instrument can record; Berkeley's ghosts are not infinitesimals
vanishing but distinctions too small to count, and the floor is exact --- a finite ruler reads in whole cells.

\section{The Precision Effect: precision is reproducibility, not truth}
\label{se:precision}

A finite history of anchor points admits a unique smoothest completion. Among all the curves that pass through
the anchors, exactly one minimizes the total curvature --- the integrated square of the second derivative,
\begin{equation}
  \int (f'')^2 ,
\end{equation}
and that minimizer is the natural cubic spline, the curve a draftsman's flexible strip traces when pinned at
the anchors, bending as little as it can between them. Minimizing the curvature forces the fourth derivative to
vanish between the anchors,
\begin{equation}
  f^{(4)} = 0,
\end{equation}
so the completion is piecewise cubic, the smoothest interpolant the anchors allow --- the completion ours to
draw, the anchors it must pass through the world's. Once the anchors are fixed
the completion is determined, and that determinacy is precision: precision is the reproducibility of the
shadow, the same smooth completion returned each time, and it is distinct from accuracy, which is nearness to a
truth the completion may or may not share. The distinction is the dartboard's. A tight cluster of throws, all
landing together but off the bull, is precise and inaccurate --- repeatable and wrong. A loose scatter centred
on the bull is accurate and imprecise --- right on average, never twice the same. An instrument can have either
without the other, and the cubic completion delivers only the first: it returns the same shadow each time, and
says nothing about whether that shadow sits on the truth.

The honest floor is a smooth-shadow analogy --- the discrete anchors are the model, the minimal-curvature
spline their smooth shadow, and the bridge from the anchors to the completion is an analogy, named as one ---
the completion ours, the truth it may or may not sit on the world's. The
reading is that precision is reproducibility, not truth. A precise instrument returns the same smooth shadow
each time it is run; whether that shadow is near the truth is accuracy, a different question, and the cubic
completion is precise exactly because the anchors fix it uniquely --- not because it is right.

\section{The Fluxions Effect: a derivative without infinitesimals}
\label{se:fluxions}

The derivative does not need a vanishing quantity. Newton's method of fluxions treated an infinitesimal $\mathrm{d}t$ as
both zero and not-zero --- nonzero while it sat in the denominator of a ratio, zero the moment the ratio was
taken --- the very paradox Berkeley mocked. But the slope can be had without it. Take the finite difference and
read off the smooth completion's slope,
\begin{equation}
  \dot x = \frac{\Delta x}{\Delta t} \longrightarrow f',
\end{equation}
the derivative being the slope of the $C^2$ cubic completion at a point, not the limit of a ratio whose
denominator disappears --- the slope ours to read, the differences that fix it the world's. No $\mathrm{d}t$
vanishes; the slope is a property of the shadow, present wherever the
completion is, needing nothing to be sent to zero.

The honest floor is a smooth-shadow analogy --- the finite differences are the model, the cubic slope their
smooth shadow, and the derivative is read off the shadow, named as such --- the shadow ours, the anchors
beneath it the world's. The reading is that the derivative is
a slope, not a vanishing. It is the slope of the smooth completion the anchors determine, and because that
completion is $C^2$ the slope is well-defined without any infinitesimal --- the fluxion's ghost laid to rest by
reading the shadow instead of vanishing the denominator.

\section{The Continuum Limit Effect: the continuum is a limit, not a substrate}
\label{se:continuum-limit}

The continuum is not a ground the discrete sits on; it is the limit the discrete approaches as the mesh
refines. A finite instrument reads on a mesh of spacing $h$, and the continuum is the $h \to 0$ limit of that
mesh --- a limit defined by the refinement study, not assumed in advance. And the limit must be earned: the
convergence law relating measurements across mesh levels is fit from the refinement, never ordained --- the
mesh ours to refine, the law its measurements obey the world's. Assume a law
in advance --- a linear stress--strain relation, say --- and the estimator sits on a plateau at $1.5$; fit it
honestly from the mesh, letting the data say how the values change as $h$ halves, and an estimator that
locates the kink converges instead to the true value, $1.0$. The number the continuum hands back depends on whether
its law was assumed or measured, and only the measured one is admissible.

The honest floor is a smooth-shadow analogy --- the mesh is the model, the continuum its smooth shadow, and the
$h \to 0$ limit is the bridge, named and earned by refinement rather than assumed --- the mesh ours, what the
refinements tend toward the world's. The reading is that the
continuum is a limit of refinements, not a substrate beneath them. It is what the mesh tends toward as it grows
fine, recovered by study and not presupposed --- the smooth not the ground the discrete stands on but the
shadow it casts as it refines.

\section{The Richardson Effect: length depends on the ruler}
\label{se:richardson}

How long is a coastline? The answer depends on the ruler. Lewis Fry Richardson, tallying the lengths of
national borders, noticed that different sources gave wildly different figures for the same frontier, and saw
why: measure Britain's coast with a finer and finer yardstick and the length does not converge but grows,
because each finer ruler resolves bays within bays, headlands within headlands, that the coarser one stepped
straight across --- the yardstick ours to pick, the bays that were always there the world's. The measured
length scales with the ruler,
\begin{equation}
  L(\epsilon) \propto \epsilon^{\,1 - D},
\end{equation}
with $D$ the fractal dimension of the coast --- a number between $1$ and $2$ measuring how rough it is. For a
smooth curve $D = 1$ and the length converges; for a coast $D > 1$ and the length diverges as $\epsilon \to 0$,
climbing without bound as the ruler shrinks. Mandelbrot made the puzzle the title of a famous paper --- how
long is the coast of Britain? --- and the answer is that there is no resolution-free length.

The honest floor is a smooth-shadow analogy --- the coastline at resolution $\epsilon$ is the model, the
continuum length its smooth shadow, and for a fractal that shadow does not exist as a single number. The
reading is that length is resolution-dependent, not absolute. A coastline has no one length, only a length at a
scale; the smooth shadow that would be ``the length'' diverges, and what is real is the ruler and the count it
returns --- the ruler ours, the count the world's --- the continuum length a shadow that, for a fractal, the
discrete never settles into.

\section{The Pythagoras--Planck Effect: the diagonal and the quantum}
\label{se:pythagoras-planck}

Here the instrument meets a no-go. A construction can force a magnitude no enumeration can reach: the diagonal
of a unit square has length $\sqrt 2$, and $\sqrt 2$ is irrational. The proof is the one that, legend says, cost its
discoverer his life --- Hippasus, the Pythagorean said to have drowned at sea for revealing that the diagonal
could not be written as a ratio, breaking the brotherhood's creed that all things were whole number. Suppose
$\sqrt 2$ were a ratio $p/q$ of whole numbers in lowest terms. Then
$p^2 = 2q^2$, so $p^2$ is even, so $p$ is even, say $p = 2k$; substituting, $4k^2 = 2q^2$, so $q^2 = 2k^2$, so
$q^2$ is even, so $q$ is even. But then $p$ and $q$ are both even, contradicting ``lowest terms.'' No ratio
survives the descent --- the ratio ours to attempt, the diagonal that refuses it the world's. So $\sqrt 2$
admits no finite ratio, no exact enumeration in whole units however fine; a
discrete ledger cannot represent it exactly, and pushed to enumerate the diagonal to the end, the refinement
never terminates. Something must give, and what gives is the assumption of unlimited refinement: the ledger
takes a smallest unit, a cap on how finely it may divide, so the irrational is approached but never exactly
reached. Physics keeps a constant at this register --- Planck's quantum of action $h$ --- and the model
borrows the name for its cap, claiming the kinship and not the derivation.

The honest floor is a no-go --- exact enumeration of the diagonal is inadmissible, and the cutoff is the
model's own, imposed on its ledger and never derived for physics --- the enumeration ours to press, the
incommensurable diagonal the world's. The reading is that the diagonal demands a smallest step of the ledger,
and the demand meets its physical echo: refinement cannot run forever against an irrational, so the ledger
takes a floor, and the name the model borrows for that floor is Planck's --- Pythagoras's diagonal and Planck's
quantum a pairing, not an identity, the kinship named and the derivation refused.

\bigskip

With the diagonal and the quantum paired, the continuum's case is closed. The continuum has been, throughout,
the most useful analogy the instrument keeps and its most carefully fenced: a resolution floor below it, a
precision that is its reproducibility, a
derivative that is its slope, a limit that earns it by refinement, a length that depends on the ruler, and an
irrational that forbids its exact enumeration. In every case the continuum is a shadow the discrete casts as it
refines --- never the substrate beneath it --- the shadow ours to name, the refinements that cast it the
world's. Part VI turns next to determinism, induction, and cause --- whether
a finite record can be continued from a seed, and what the gap between past and future costs.

\bigskip
\begin{center}
\rule{0.4\textwidth}{0.4pt}
\end{center}
\bigskip

\noindent The case asks, before the end, how finely its readings can be made --- how sharply the world can be told
apart --- and in the asking it meets the one thing it has held at arm's length the whole book long: the
continuum, the unbroken line physics is written on. It is the instrument's most useful companion and the one it
has handled most warily, for the continuum is never a thing the instrument holds; it is a shadow the instrument
casts.

Every reading has a floor. Below some smallest step the instrument records nothing --- two states nearer than
that are, to it, one state, and no care sharpens what lies beneath the floor; a finer ruler only lowers the
floor, it does not abolish it. This is the resolution the very first chapter laid down: the grid we set, below
which nothing varies. Galileo, grinding lenses to push the floor down, saw moons around Jupiter the naked eye
could not --- not because the moons were tiny beyond measure but because they lay below the eye's floor and
above the glass's; and Berkeley's mockery of the infinitesimal, the "ghosts of departed quantities," is answered
by exactly this floor: there are no departed quantities, only distinctions too small to count. And what the instrument returns above the floor is precise when it returns the same shadow
each time --- but precision is not truth: a tight cluster of readings can sit, repeatably, off the mark, and
whether they are right is a separate question the sameness never answers. The smooth curve the instrument draws
through its readings is just that --- a curve drawn \emph{through}, the shape the discrete anchors approach, and
not a thing lying beneath them. Even how long a rough thing is depends on the ruler that measures it: shrink the
yardstick and a coastline's length climbs without end, for there was never one length, only a length at a scale
--- how long is the coast of Britain, the famous question runs, and the honest answer is that there is no
resolution-free length to give.

So the continuum stands where it has stood since the first count: the smooth shadow the discrete casts as it is
refined, never the ground the discrete rests on. It is the most useful analogy the instrument keeps --- the
bridge to every equation, drawn again and again --- and the most carefully fenced, offered always as a
convenience of ours and never as a fact the world hands over. Press the ledger to write out an irrational ---
the diagonal of a plain square, which no ratio of whole numbers can ever equal --- and the refinement never
finishes: the ledger cannot reach it, only approach, and must at the last take a smallest step and stop. Legend says the
first man to prove that diagonal beyond the reach of any ratio drowned for telling of it, so hard did the old
creed that all things are whole numbers die. The
instrument names the continuum, uses it, leans its whole calculus on it --- and marks, each time, that it is a
shadow named as a shadow, a limit approached and not a floor to stand on. What more there is to say of it, the
case keeps for the end.

And the \emph{charge}, read this finely, meets a floor of its own. The count the case will collect on is read to
a smallest step and no further; below it two readings are one, and no sharpening tells them apart. The charge is
exact to the floor and silent beneath it --- a count in whole cells, not a real number carried out forever. It
is never owed to more places than the reading has; a charge claimed sharper than the ruler that took it would be
a precision the record does not hold, and the instrument will not pretend to it. And
the continuum that would smooth the count into an unbroken quantity is, here as everywhere, a shadow of the
count and not its ground: the tally is the real thing, the smooth curve through it a convenience. The trace,
likewise, is a count kept to its floor --- what the record retains, retained in whole steps, exact as far as it
reads and honest about where it stops. Whatever comes to be owed is owed in counted cells, to the floor of the
reading and no finer.

And each of the three readings the case will bring has a floor and a precision of its own --- its own smallest
step, its own reproducibility. To meet on one number the three must agree not to a perfection none of them holds
but to the coarsest floor among them; agreement finer than the roughest ruler would be a fiction. The three are
corrected and ready, each honest about how finely it reads; the number they must meet on is still not in hand. The
floor of the reading is one more thing the meeting will have to respect.

The instrument has one inward question left before the reckoning. Having asked how finely a record can be read, it
asks next whether a record can be carried forward at all --- whether a finite past fixes its future or leaves a
gap between them, and what that gap, if it is there, costs. The next chapter turns to determinism, induction,
and cause. The case is still open, its readings admissible and floored and not yet agreed; the court does not
sit until the record is complete.
